\section{Checks, provenance and mathematical status}\label{app:checks}
The analytical arguments in this draft are written proofs and explicitly
qualified construction problems. Eight small standard-library Python
programs accompany the investigation. They check finite algebra and
normalization identities rather than positivity of the Weil form.
The preserved records distinguish exact rational checks from labelled
floating-point illustrations.
\begin{center}
\small
\begin{tabular}{p{0.38\linewidth}rp{0.40\linewidth}}
\toprule
Program in \texttt{numerics/} & Cases & Scope \\
\midrule
\texttt{check\_matching.py} & 153 & Rational filters and gamma refinement \\
\texttt{check\_gauge\_transfer.py} & 326 & Winding, transfer and first-prime algebra \\
\texttt{check\_sphere\_schur.py} & 251 & Charged norms, product ratios and undressed overlaps \\
\texttt{check\_dressed\_schur.py} & 587 & Dressed norm, magnetic coefficients and mixed-return witnesses \\
\texttt{check\_explicit\_formula.py} & 8 & Normalization against the sum over zeta zeros \\
\texttt{check\_channel\_bound.py} & 75 & Interference bound and the two-channel certificates \\
\texttt{check\_mirror\_subtraction.py} & 75 & Mirror weight, subtraction identity, domination quotients \\
\texttt{check\_prime\_free\_archimedean.py} & 78,666 & Parity split, pole factorisation, the criterion of Proposition~\ref{prop:prime-free-large-L} at every prime breakpoint below $10^6$ \\
\bottomrule
\end{tabular}
\end{center}
The 80,141 finite cases are supporting checks, not 80,141 independent
theorems. Finite-product end factors and truncation boundary
coefficients are retained. The finite mixed-return sample does not
replace the unique-factorization proof in Proposition~\ref{prop:mixed}.
The Mellin and beta integrals, infinite-series convergence,
graph-domain argument and kernel exclusions are justified in the text;
the programs do not formally verify them.

The explicit-formula, channel-bound, mirror-subtraction and prime-free
programs are labelled floating-point computations rather than exact rational
ones. \texttt{check\_explicit\_formula.py} evaluates
both sides of \eqref{eq:explicit-formula} for the input
$F(x)=\tfrac12(1+\cos(2\pi x/L))$ on $I_L$, whose transform is
$\widehat F(\tau)=\omega^2\sin(\tau L/2)/(\tau(\omega^2-\tau^2))$ with
$\omega=2\pi/L$. It computes the zero ordinates itself, from the Hardy
function through Borwein's series for $\eta$, and agrees with published
values to $10^{-13}$. For $L=2,3,4$ the two sides agree to
$2.5\times10^{-10}$, $5.4\times10^{-11}$ and $1.0\times10^{-11}$, against
totals of $2.9\times10^{-5}$, $2.6\times10^{-6}$ and $1.4\times10^{-7}$:
the ratio of the largest single term of \eqref{eq:target} to the total rises
from $1.4\times10^5$ to $6.3\times10^7$ across those three intervals. This
check confirms the conventions of \eqref{eq:weil-functional} as a whole. It
is not evidence for RH, since its ordinates are computed on the critical
line by construction.
\texttt{check\_channel\_bound.py} evaluates \eqref{eq:two-channel-test} for
step-function inputs, for which every quantity is a finite closed form; the
one infinite series, $\sum_n(\cdot)e^{-a_nr}$, is summed after subtracting
its exactly known leading part, and the reported remainder bound for the two
certificates of Theorem~\ref{thm:two-channel} is below $10^{-40}$.
\texttt{check\_mirror\_subtraction.py} verifies
\eqref{eq:mirror-coefficients} and \eqref{eq:complementary} in exact rational
arithmetic, checks \eqref{eq:subtraction} as an identity of quadratic forms on
step inputs, and reports the domination quotients discussed after
Corollary~\ref{cor:domination}. Those quotients are one sided: a value above
one on a subspace certifies that a domination fails, a value below one
certifies nothing. The Galerkin eigenvalues quoted in
Section~\ref{subsec:subtraction} are labelled floating-point computations and
are not part of the replayed record.

\texttt{check\_prime\_free\_archimedean.py} builds the forms of
\eqref{eq:target} and \eqref{eq:prime-free} in a uniform cell basis, checks
that basis against the independent step routine of the channel-bound program,
verifies \eqref{eq:parity-split} and \eqref{eq:hyperbolic} on random step
inputs, and evaluates \eqref{eq:large-L-criterion} at all $78{,}497$ prime
breakpoints below $10^6$. Its Galerkin quotients are one sided in the same
sense as the mirror-subtraction ones. Two of its outputs are conclusive rather
than indicative, because a negative value in a subspace certifies that a form
is not positive: the indicator of $I_L$ gives
$K[E_L1]+w_0\norm1^2<0$ for $L\leq3/2$, so the pole term of
\eqref{eq:prime-free-target} cannot be dispensed with; and an explicit step
function on $I_{3/4}$ gives $K+w_0\norm\cdot^2+P_L$ the value
$-3.2\times10^{-3}\norm f^2$, so the archimedean form with the bare constant
$w_0$ and no mirror contact is not positive for any $L\geq3/4$. Since
$\widetilde\kappa_{r_2,d_2}=0.574$ switches on at $L=\log2$, the mirror
contact of Theorem~\ref{thm:subtraction} is what keeps \eqref{eq:prime-free}
positive beyond that point.

The manuscript is self-contained at the level of its stated external
inputs: the classical Weil criterion and the cited model representation
data. It does not input a shared background file or cite equation numbers
from one. The accompanying coverage map relates its sections to the
retained research notes, including the broader candidate survey.
Build instructions and a hash record identify the reviewed sources and
compiled PDF. The \texttt{drafts/} folder preserves complete dated
source-and-PDF snapshots. Those hashes establish file identity, not
mathematical correctness or portability of PDF metadata between TeX
installations.

The following boundaries of the results are essential. The positive
prime reference is known; the dressed theorem gives an explicit
elementary Schur source for it. Square-summable formal magnetic
outputs do not settle all closed-operator extensions. The continuous
lift has not been derived from a four-dimensional interface. The
general protected-sector construction hypotheses have not been
proved here. Theorem~\ref{thm:two-channel} excludes one specified
architecture and not a class of theories. Most importantly, the full
contact-and-pole identity in Problem~\ref{prob:joint} remains unsolved, so
this draft makes no claim to establish RH or a completed reduction to
constructive QFT.
